\chapter{Kilder, opdatering og videre læsning}

\chapterlead{En juridisk bog er aldrig en erstatning for den gældende retskilde. Brug dette kapitel som et kort over, hvor du kontrollerer lovtekst, praksis, vejledning og ændringer.}

\section{Officielle kilder}

\small
\begin{longtable}{@{}p{3.1cm}@{}p{9.9cm}@{}}
\toprule
\textbf{Kilde} & \textbf{Hvad finder du her?} \\
\midrule
Retsinformation & Gældende danske love, bekendtgørelser, cirkulærer, lovforslag og dokumenter. \url{https://www.retsinformation.dk/} \\
Danmarks Domstole & Information om domstole, retssager, sagsforløb og praktiske procedurer. \url{https://domstol.dk/} \\
EUR-Lex & EU-traktater, forordninger, direktiver, afgørelser og EU-Domstolens materiale. \url{https://eur-lex.europa.eu/} \\
ECHR & Den Europæiske Menneskerettighedskonvention, protokoller og materiale om domstolen. \url{https://www.echr.coe.int/european-convention-on-human-rights} \\
Datatilsynet & Dansk vejledning og tilsynsmateriale om databeskyttelse. \url{https://www.datatilsynet.dk/regler-og-vejledning/lovgivning} \\
Folketinget & Lovforslag, debatter, udvalg og lovgivningsmateriale. \url{https://www.ft.dk/} \\
Folketingets Ombudsmand & Tilsyn med offentlig forvaltning og udtalelser om god forvaltningsskik. \url{https://www.ombudsmanden.dk/} \\
\bottomrule
\end{longtable}
\normalsize

\section{Centrale tekster}

I bogens kapitler er de vigtigste officielle tekster henvist gennem bibliografien. De omfatter Grundloven, forvaltningsloven, offentlighedsloven, retsplejeloven, straffeloven, aftaleloven, købeloven, erstatningsansvarsloven, databeskyttelsesloven, GDPR, EU-charteret, AI-forordningen og ECHR \citep{grundloven,forvaltningsloven,offentlighedsloven,retsplejeloven,straffeloven,aftaleloven,koebeloven,erstatningsansvar,databeskyttelsesloven,gdpr,ai-act,charter,echr}.

\section{Søgning efter ændringer}

Lovtekster ændres. En ny ændringslov kan være vedtaget uden at være trådt i kraft, og en konsolideret tekst kan indeholde senere ændringer med forskellige ikrafttrædelsesdatoer. Kontroller derfor:
\begin{enumerate}
  \item dokumentets status: gældende, historisk eller foreslået,
  \item seneste ændringer,
  \item ikrafttrædelsesbestemmelsen,
  \item overgangsregler, og
  \item om sagen ligger i Danmark, på Færøerne, i Grønland eller i en anden jurisdiktion.
\end{enumerate}

Den samme lov kan have særlige regler om geografisk anvendelse. En tekst, der gælder i Danmark, gælder ikke automatisk i alle dele af rigsfællesskabet eller i andre lande.

\section{Kildekritik}

En blog, et forumindlæg eller en AI-genereret tekst kan være et nyttigt sted at opdage et søgeord. Det er ikke nødvendigvis en autoritativ kilde til gældende ret. Kildekritik handler om:
\begin{itemize}
  \item hvem der har udstedt eller skrevet materialet,
  \item om teksten er aktuel,
  \item om den beskriver lov, praksis eller holdning,
  \item om den angiver sin jurisdiktion, og
  \item om konklusionen kan kontrolleres i en primær kilde.
\end{itemize}

\section{Versionsnotat}

Denne udgave er skrevet til et dansk introduktionsniveau og kontrolleret med officielle kilder den 15. juli 2026. Den er bevidst skrevet med generelle forklaringer frem for en udtømmende paragrafkommentar. Ved senere brug i praksis eller videre læsning bør læseren kontrollere senere ændringer og supplere med aktuelle domme, vejledninger og speciallitteratur.

\caution{Hvis du står med en konkret frist, risiko for straf, udsættelse, tvangsfuldbyrdelse, alvorlig familiesag eller stort økonomisk krav, så brug ikke denne bog som eneste grundlag for handling. Find den gældende regel og søg kvalificeret hjælp.}
